\documentclass[11pt]{article}
\usepackage[margin=1.1in]{geometry}
\usepackage{amsmath,amssymb,amsthm}
\usepackage{booktabs}
\usepackage[round]{natbib}
\usepackage[hidelinks]{hyperref}
\hypersetup{pdftitle={Gaussian Perturbed Maxima Without Sampling},
pdfauthor={Peter Cotton}}

\newtheorem{proposition}{Proposition}
\theoremstyle{remark}
\newtheorem{remark}{Remark}

\newcommand{\E}{\mathbb{E}}
\newcommand{\Prob}{\mathbb{P}}

\title{Gaussian Perturbed Maxima Without Sampling}
\author{Peter Cotton\thanks{Working draft. Every identity below is
verified numerically (experiment 27 of the companion repository,
\url{https://github.com/microprediction/kinetics}); the neural
experiments of Section~\ref{sec:program} are pending.}}
\date{August 2026}

\begin{document}
\maketitle

\begin{abstract}
Softmax is the gradient of a perturbed maximum: independent Gumbel noise
turns the multiclass maximum into log-sum-exp. General perturbations
give other smooth choice maps and losses, but their expectations are
evaluated by Monte Carlo, and the leading heteroscedastic image
classifiers, whose output layer is exactly a low-rank Gaussian perturbed
argmax, approximate their probabilities by sampling a softmax
relaxation. For Gaussian perturbations with low-rank-plus-diagonal
covariance we make the whole object explicit and sampling-free. One
shared conditional field yields every winner probability, the potential,
exact Jacobian--vector products, and, by a new reverse-mode formula, the
gradients with respect to means, factor loadings, and idiosyncratic
variances, all at cost linear in the number of alternatives. The dual
regularizer's curvature is the effective resistance of a photo-finish
graph, with softmax the special case whose conductances are $p_i p_j$.
The likelihood and Fenchel--Young objectives separate outside that
special case, and learning covariance by the latter is degenerate.
\end{abstract}

\noindent\textbf{Keywords:} perturbed optimizers; multinomial probit;
Fenchel--Young losses; graph Laplacian; heteroscedastic classification;
effective resistance

\section{Introduction}

A large class of neural operations is organized around the maximum of
$N$ competing scores. The maximizer is discrete, so it is smoothed by
perturbation: for noise $\xi$, define
\[
G(\mu) = \E\Big[\max_i (\mu_i + \xi_i)\Big], \qquad
p(\mu) = \nabla G(\mu) = \E\big[e_{\arg\max_i(\mu_i + \xi_i)}\big].
\]
For independent Gumbel noise $G$ is log-sum-exp and $p$ is softmax, up
to an additive constant. \citet{berthet2020} develop this perturbed-optimizer view in
generality, including the convex dual and the induced Fenchel--Young
loss \citep{blondel2020}, and evaluate general perturbations by Monte
Carlo; they demonstrate Gaussian-perturbed argmax training on CIFAR-10.

The Gaussian member is the natural carrier of correlation. We study
\[
U_i = \mu_i + v_i^\top f + \sqrt{D_i}\,\epsilon_i, \qquad
f \sim N(0, I_k), \quad \epsilon_i \sim N(0,1) \text{ iid},
\]
so the perturbation covariance is $\Sigma = VV^\top + \mathrm{diag}(D)$,
low-rank plus diagonal. This structure arises automatically from
uncertain representations: if a feature vector is $h = \bar h + Af$ and
scores are $U_i = w_i^\top h + b_i + \sqrt{D_i}\,\epsilon_i$, then
$v_i = A^\top w_i$.

The same model already ships in computer vision.
\citet{collier2021} place exactly this latent-utility argmax, with
input-dependent low-rank-plus-diagonal covariance, on the final layer of
ImageNet-scale classifiers; lacking a closed form, they approximate the
winner probabilities by Monte Carlo over a temperature-softmax
relaxation, and the scaled follow-up \citep{collier2023} identifies the
latency and memory cost of that sampling as a standing issue.

A numerical companion \citep{cotton2026} supplies the missing
primitive. Conditional on the factors the alternatives are independent,
so all $N$ winner probabilities come from a single product field in
$O(QNL)$ operations for $Q$ factor nodes and $L$ lattice points, with
exact Jacobian--vector products from the same pass, a weighted
graph-Laplacian Jacobian, and a globally invertible map from mean
utilities modulo translation to the simplex interior.

This paper asks what that primitive becomes as a neural layer. The
novelty boundary is deliberately narrow. Perturbed-max smoothing, its
duality, and Fenchel--Young losses are prior art
\citep{berthet2020,blondel2020}; generalized-entropy conjugates of
random-utility surpluses are prior art \citep{fosgerau2020}. What we add
is what becomes computationally and geometrically explicit in the
structured Gaussian case:

\begin{enumerate}
\item \textbf{The potential is free.} The field that computes $p$
computes $G$ with $O(QL)$ additional work
(Section~\ref{sec:potential}).
\item \textbf{Dual graph geometry.} The dual regularizer's Hessian is
the pseudoinverse of the photo-finish Laplacian, so dual curvature is
effective electrical resistance; softmax is the special graph with
conductances $p_i p_j$ (Section~\ref{sec:dual}).
\item \textbf{Full reverse mode.} A new shared-field vector--Jacobian
product gives gradients in $(\mu, V, D)$ at the cost of one forward
pass (Section~\ref{sec:reverse}).
\item \textbf{Two losses, one special case.} The Gaussian
Fenchel--Young gradient is $p - e_y$; the exact likelihood gradient is
$-J e_y / p_y$, which routes the label through its photo-finish edges.
They coincide exactly when the edge weights are $p_i p_j$, i.e.\ for
softmax (Section~\ref{sec:losses}).
\item \textbf{A covariance-learning caution.} The heat identity
$\partial G/\partial D_i = J_{ii}/2 > 0$ and $\nabla_V G = JV$ make the
Fenchel--Young loss monotone in choice-relevant perturbation variance:
it cannot identify free Gaussian noise (Section~\ref{sec:losses}).
\end{enumerate}

Every displayed identity is verified numerically against finite
differences, Monte Carlo, or direct linear algebra
(Table~\ref{tab:identities}); the neural experiments of
Section~\ref{sec:program} are pending.

\section{The structured Gaussian perturbed maximum}

Conditional on $f$, the scores are independent with means
$m_i(f) = \mu_i + v_i^\top f$ and variances $D_i$. Write $F_i(x \mid f)$
and $g_i(x \mid f)$ for the conditional CDF and density. Then
\[
p_i = \Prob(U_i = \max_j U_j)
= \E_f \int g_i(x \mid f) \prod_{j \ne i} F_j(x \mid f)\, dx .
\]
The shared field is $H_f(x) = \prod_j F_j(x \mid f)$; each winner
integrand divides one factor back out. On $Q$ factor nodes and an $L$-point
lattice, all $N$ probabilities cost $O(QNL)$ \citep{cotton2026}.

\section{The potential from the same field}\label{sec:potential}

Conditional on $f$, $H_f$ is the CDF of the maximum, so for the lattice
endpoint $b$,
\[
\E\big[\max_i U_i \,\big|\, f\big] = b - \int_{-\infty}^{b} H_f(x)\, dx ,
\]
and $G(\mu)$ is the factor average. The forward pass has already built
$H_f$ at every node and lattice point; the potential is one more
reduction.

One discretization point is not cosmetic. The continuum identity
$\nabla G = p$ survives the lattice only if the potential's quadrature
is consistent with the probability quadrature: a rectangle-rule
potential next to the implemented probability integrals broke the
identity at $2.5\times10^{-4}$ in our tests, while the trapezoidal
form restores it to $3.6\times10^{-8}$. Fenchel--Young training needs a
potential-consistent mode, with the discrete $\nabla G = p$ defect
reported, not assumed away.

\section{Dual geometry}\label{sec:dual}

$G$ is convex and $G(\mu + c\mathbf{1}) = G(\mu) + c$, so the conjugate
$\Psi = G^*$ lives on the simplex and, on the mean-zero gauge,
$p = \nabla G(\mu)$ and $\mu = \nabla \Psi(p)$: the share inversion of
the companion paper is the mirror map, and
$\Psi(p) = p^\top \mu(p) - G(\mu(p))$ is computable from the inverse
transform plus the same-pass potential.

The Hessian is the photo-finish Laplacian
\[
J = \nabla^2 G = \sum_{i<j} w_{ij} (e_i - e_j)(e_i - e_j)^\top,
\qquad
w_{ij} = \E_f \int g_i g_j \prod_{\ell \ne i,j} F_\ell \, dx > 0,
\]
with $w_{ij}$ the density of an $i$--$j$ photo finish for the winning
position. On the tangent space, $\nabla^2 \Psi(p) = J^+$. For the
transfer direction $d_{ij} = e_i - e_j$,
\[
d_{ij}^\top J^+ d_{ij} = \text{effective resistance between $i$ and $j$}
\]
in the network with conductances $w_{ab}$: probability moves cheaply
between strongly competing alternatives and expensively across weak
competitive pathways.

Softmax is the special case $w_{ij} = p_i p_j$, whose resistance is
$1/p_i + 1/p_j$ --- exactly the curvature of negative Shannon entropy
along $d_{ij}$. Shannon geometry is the effective-resistance geometry of
one particular competition graph. A structured Gaussian layer separates
information softmax cannot: two states with identical $p$ can carry
different photo-finish graphs, distinguishing ``similar marginal
probability'' from ``direct competitors at the decision boundary''.

\section{Shared-field reverse mode}\label{sec:reverse}

The companion paper differentiates $p$ in $\mu$. A heteroscedastic layer
needs $V$ and $D$ as well. Let $r_i(x, f) = g_i \prod_{\ell \ne i}
F_\ell$ be the winner densities, let a downstream loss supply the
cotangent $c_i = \partial \mathcal{L}/\partial p_i$, and form one
additional field $S_c = \sum_i c_i r_i$ with hazards
$\lambda_j = g_j / F_j$.

\begin{proposition}[Shared-field vector--Jacobian product]
\label{prop:vjp}
With $a_j(f) = \int \big[ c_j r_j \frac{x - m_j}{D_j} - (S_c - c_j r_j)
\lambda_j \big] dx$ and $b_j(f) = \int \big[ c_j r_j
\frac{(x-m_j)^2 - D_j}{2 D_j^2} - (S_c - c_j r_j)
\frac{x - m_j}{2 D_j} \lambda_j \big] dx$,
\[
\frac{\partial \mathcal{L}}{\partial \mu_j} = \E_f[a_j], \qquad
\frac{\partial \mathcal{L}}{\partial v_j} = \E_f[a_j f], \qquad
\frac{\partial \mathcal{L}}{\partial D_j} = \E_f[b_j].
\]
All $N$ parameter blocks reuse the same $S_c$; the cost is
$O(QN(L + k))$.
\end{proposition}

\begin{proof}
Only two terms depend on alternative $j$'s parameters: its own winner
density through $g_j$, and every other winner density through $F_j$.
Hence $\partial_{\theta_j} \mathcal{L} = \E_f \int [ c_j r_j\,
\partial_{\theta_j} \log g_j + (S_c - c_j r_j)\, \partial_{\theta_j}
\log F_j ]\, dx$. For a Gaussian factor, $\partial_{m_j} \log g_j =
(x - m_j)/D_j$ and $\partial_{m_j} \log F_j = -\lambda_j$;
$\partial_{D_j} \log g_j = ((x-m_j)^2 - D_j)/(2 D_j^2)$ and
$\partial_{D_j} \log F_j = -\lambda_j (x - m_j)/(2 D_j)$. Finally
$\partial_{\mu_j} m_j = 1$ and $\partial_{v_j} m_j = f$.
\end{proof}

Measured against central finite differences (experiment 27), the three
blocks agree to $2.2\times10^{-9}$, $4.9\times10^{-10}$, and
$6.5\times10^{-10}$.

If the implementation normalizes, $p = \tilde p / s$ with
$s = \mathbf{1}^\top \tilde p$, the incoming cotangent must first be
pulled back: $\tilde c = (c - (c^\top p)\mathbf{1})/s$. Differentiating
the unnormalized field while treating normalization as constant is
incorrect.

\section{Two learning objectives}\label{sec:losses}

\paragraph{Fenchel--Young.} For a one-hot label $e_y$ and centered
Gaussian noise, $\Psi(e_y) = 0$, so the Fenchel--Young loss induced by
$G$ is
\[
L_{\mathrm{FY}}(\mu, y) = G(\mu) - \mu_y, \qquad
\nabla_\mu L_{\mathrm{FY}} = p - e_y ,
\]
the exact Gaussian counterpart of the softmax cross-entropy gradient.
In the binary case, with $\delta = \mu_1 - \mu_2$ and
$s^2 = \mathrm{Var}(\xi_1 - \xi_2)$,
$L_{\mathrm{FY}} = s[\phi(\delta/s) - (\delta/s)\Phi(-\delta/s)] =
\E[(U_2 - U_1)_+]$: a Gaussian-smoothed hinge whose training signal
concentrates on photo finishes (verified to $4.9\times10^{-13}$
against the field computation).

\paragraph{Likelihood.} Treating the race as a probabilistic model
gives $L_{\mathrm{NLL}} = -\log p_y$, and by symmetry of $J$,
\[
\nabla_\mu L_{\mathrm{NLL}} = -\frac{J e_y}{p_y}
= \frac{1}{p_y} \sum_{j \ne y} w_{yj} (e_j - e_y):
\]
the label updates only through its photo-finish edges. For softmax,
$w_{yj} = p_y p_j$ collapses this to $p - e_y$: the two objectives
coincide exactly in the Gumbel/softmax case and separate otherwise.
Their curvatures likewise separate: Fenchel--Young curvature is $J$;
the Fisher information of the winner likelihood is
$J\,\mathrm{diag}(p)^{-1} J$, and softmax satisfies
$J\,\mathrm{diag}(p)^{-1} J = J$.

\paragraph{Covariance collapse.} The Gaussian potential obeys the heat
identity $dG = \tfrac12 \mathrm{tr}(J\, d\Sigma)$, hence
\[
\frac{\partial G}{\partial D_i} = \frac{J_{ii}}{2} > 0, \qquad
\nabla_V G = J V
\]
(both verified to $\sim10^{-8}$). Since $-\mu_y$ carries no covariance
dependence, $\nabla_{V,D} L_{\mathrm{FY}} = \nabla_{V,D} G$: the
Fenchel--Young loss is monotone increasing in every positive-semidefinite
direction of choice-relevant perturbation covariance. Unconstrained
minimization removes the noise. If $\Sigma$ is a chosen smoothing
geometry, fix or constrain it; if $\Sigma(x)$ is uncertainty to be
learned, use the winner likelihood or an outer criterion.

\paragraph{Gauges.} An end-to-end layer must remove three invariances:
$\mu \mapsto \mu + c\mathbf{1}$; $V \mapsto V + \mathbf{1}a^\top$
(common factor noise); and joint scaling $(\mu, V, D) \mapsto
(s\mu, sV, s^2 D)$. Project $\mu$ and $V$ onto mean-zero columns and
normalize the contrast covariance, e.g.\
$\mathrm{tr}(P \Sigma P)/(N-1) = 1$ with
$P = I - \mathbf{1}\mathbf{1}^\top/N$. Factor rotation
$V \mapsto VR$ remains unidentified and is harmless for prediction.

\begin{table}[t]
\centering
\begin{tabular}{lr}
\toprule
identity & max error \\
\midrule
$\nabla_\mu G = p$ & $3.6\times10^{-8}$ \\
reverse mode, $\mu$ block & $2.2\times10^{-9}$ \\
reverse mode, $V$ block & $4.9\times10^{-10}$ \\
reverse mode, $D$ block & $6.5\times10^{-10}$ \\
heat identity $\partial G/\partial D_i = J_{ii}/2$ & $1.0\times10^{-8}$ \\
$\nabla_V G = JV$ & $7.7\times10^{-9}$ \\
$\nabla_\mu(-\log p_y) = -J e_y / p_y$ & $5.1\times10^{-9}$ \\
binary Fenchel--Young closed form & $4.9\times10^{-13}$ \\
dual curvature $=$ effective resistance & $1.4\times10^{-6}$ (rel.) \\
softmax $R_{ij} = 1/p_i + 1/p_j$ & $1.8\times10^{-15}$ \\
$G$ against $4\times10^6$-draw Monte Carlo & $1.9\times10^{-4}$ ($\tfrac12$ s.e.) \\
\bottomrule
\end{tabular}
\caption{Every claimed identity, verified (experiment 27: $N=5$, $k=1$,
Gauss--Hermite 21, $L=2001$, against central finite differences, Monte
Carlo, or direct linear algebra).}
\label{tab:identities}
\end{table}

\section{Experimental program}\label{sec:program}

Three questions, kept separate. \textbf{A} (numerical): can the
operator be evaluated and differentiated accurately without resampling?
\textbf{B} (optimization): does a fixed accurate gradient beat $M$-sample
Monte Carlo gradients at matched wall time --- including $M = 1$, which
\citet{berthet2020} note is often adequate for SGD? \textbf{C}
(statistical): does structured Gaussian competition improve the model?
A faster operator is not evidence of a better inductive bias, and
vice versa.

The planned sequence: (I) operator microbenchmarks over
$N \in \{10, \dots, 10^4\}$, $k \le 15$ (product Gauss--Hermite at low
$k$, fixed scrambled Sobol above), reporting probability, potential,
and gradient error against wall time and memory, with equal-time and
equal-error comparisons against hard-argmax Monte Carlo, common-random-number
and RQMC variants, and the temperature-softmax surrogate; (II) the
identity suite of Table~\ref{tab:identities} extended to automated
randomized tests plus autograd \texttt{gradcheck}; (III) frozen-feature
classification (convex in the head) comparing softmax, Gaussian
Fenchel--Young at $M \in \{1, \dots, 1000\}$, the resampling-free
evaluator, and exact winner NLL, on objective gap and accuracy versus
wall clock; (IV) replication of the Gaussian CIFAR-10 experiment of
\citet{berthet2020} with one added treatment, the deterministic
evaluator; (V) structured fixed covariance from held-out confusion,
prototypes, or augmentation logit covariance at $k \in \{1,2,4,8\}$;
(VI) the heteroscedastic classifier of \citet{collier2021} re-trained
with the exact winner likelihood via Proposition~\ref{prop:vjp},
against the Monte-Carlo temperature-softmax surrogate, both evaluated
under a common high-accuracy reference.

Success requires any of: materially better accuracy-per-wall-time than
Monte Carlo; materially lower gradient error at matched time; faster
convergence to a validation threshold; large latency or memory savings
at high accuracy; or better-calibrated models from the exact
likelihood. The project pivots if $M=1$ Monte Carlo trains as well and
faster, if useful ranks are unreachable, or if structure adds no
statistical gain --- in which case the operator's role is high-accuracy
inference and calibration, not SGD replacement.

\section{Discussion}

Softmax is the explicitly computable Gumbel member of the perturbed-max
family. The results above make the low-rank Gaussian member explicit
too: potential, probabilities, inverse, and all parameter gradients
from one shared field, with a graph geometry that softmax collapses to
a single special case. The claim is not that Gaussian perturbations
dominate Gumbel ones; it is that in this structured regime,
computability no longer decides the comparison in advance. The
heteroscedastic classification literature already states the demand:
its generative model is this model, and its current evaluator is a
sampled relaxation of it.

\section*{Reproducibility}

Experiment 27 in the companion repository implements the operator and
verifies every identity in Table~\ref{tab:identities}; the companion
paper's experiments 13--26 cover the underlying transform, its
certification, and its benchmarks.

\begin{thebibliography}{9}

\bibitem[Berthet et~al.(2020)Berthet, Blondel, Teboul, Cuturi, Vert,
and Bach]{berthet2020} Q.~Berthet, M.~Blondel, O.~Teboul, M.~Cuturi,
J.-P.~Vert, F.~Bach. Learning with differentiable perturbed optimizers.
\emph{Advances in Neural Information Processing Systems} 33 (2020)
9508--9519.

\bibitem[Blondel et~al.(2020)Blondel, Martins, and
Niculae]{blondel2020} M.~Blondel, A.~F.~T. Martins, V.~Niculae.
Learning with Fenchel--Young losses. \emph{Journal of Machine Learning
Research} 21 (2020) 1--69.

\bibitem[Collier et~al.(2021)Collier, Mustafa, Kokiopoulou, Jenatton,
and Berent]{collier2021} M.~Collier, B.~Mustafa, E.~Kokiopoulou,
R.~Jenatton, J.~Berent. Correlated input-dependent label noise in
large-scale image classification. \emph{Proceedings of CVPR} (2021)
1551--1560.

\bibitem[Collier et~al.(2023)Collier, Jenatton, Mustafa, Houlsby,
Berent, and Kokiopoulou]{collier2023} M.~Collier, R.~Jenatton,
B.~Mustafa, N.~Houlsby, J.~Berent, E.~Kokiopoulou. Massively scaling
heteroscedastic classifiers. \emph{Proceedings of ICLR} (2023).

\bibitem[Cotton(2021)]{cotton2021} P.~Cotton. Inferring relative
ability from winning probability in multientrant contests. \emph{SIAM
Journal on Financial Mathematics} 12 (2021) 295--317.

\bibitem[Cotton(2026)]{cotton2026} P.~Cotton. Scalable probit share
calibration. Working paper (2026),
\url{https://github.com/microprediction/kinetics}.

\bibitem[Fosgerau et~al.(2020)Fosgerau, Melo, de~Palma, and
Shum]{fosgerau2020} M.~Fosgerau, E.~Melo, A.~de~Palma, M.~Shum.
Discrete choice and rational inattention: a general equivalence result.
\emph{International Economic Review} 61 (2020) 1569--1589.

\end{thebibliography}

\end{document}
